\section*{Capítulo \ref{ch:b2:field-lab-methods-2} --- Métodos de campo y de laboratorio II}
% numbers from the Book 4 second-writer scratch calculator calc/ch31.py

\begin{solution}{exo:b2:field-lab-methods-2:1}
(a) (i) \omterm{def:b2:field-lab-methods-2:textures}{Textura ofítica}; (ii) \omterm{def:b2:field-lab-methods-2:textures}{textura granoblástica}; (iii) un \omterm{def:b2:field-lab-methods-2:porphyroblast}{poiquiloblasto}, un \omterm{def:b2:field-lab-methods-2:porphyroblast}{porfiroblasto} lleno de
inclusiones; (iv) una \omterm{def:b2:field-lab-methods-2:porphyroblast}{corona}, un borde de reacción alrededor de un granate que se estaba descomponiendo. (b) Por
\eqref{eq:b2:field-lab-methods-2:lami}, $\gamma_1/\gamma_2 = \sin\ang{104}/\sin\ang{128} =
0.9703/0.7880 = 1.23$: ese límite tira un \qty{23}{\percent} más fuerte que los otros dos.
\end{solution}

\begin{solution}{exo:b2:field-lab-methods-2:2}
(a) $Q + A + P = 80$: Q $= 100 \times 25/80 = 31.3$, A $= 18.8$, P $= 50.0$. (b) P/(A+P) $=
40/55 = \qty{72.7}{\percent}$, entre 65 y 90 con Q entre 20 y 60: una granodiorita.
(c) $M' = 12 + 8 = \qty{20}{\percent}$, por debajo de 35: leucocrática.
\end{solution}

\begin{solution}{exo:b2:field-lab-methods-2:3}
(a) $\sigma = \sqrt{0.12 \times 0.88/800} = 0.0115$, es decir, $12.0 \pm 1.1$ puntos porcentuales, un
error relativo de $0.0115/0.12 = \qty{9.6}{\percent}$. (b) Cuatro veces más puntos lo reducen a la mitad: $0.0057$.
\end{solution}

\begin{solution}{exo:b2:field-lab-methods-2:4}
(a) La pérdida es de \qty{0.0288}{g}, y $0.0288/0.8000 = \qty{3.60}{\percent}$.
(b) $0.7712/8.7712 = 0.0879$.
(c) Polvo calcinado: $5.20/0.08792 = \qty{59.1}{\percent}$; polvo original:
$59.14 \times (1 - 0.0360) = \qty{57.0}{\percent}$.
\end{solution}

\begin{solution}{exo:b2:field-lab-methods-2:5}
(a) $\sigma = 0.10 \times 0.02 = 0.002$ y $n = 0.02 \times 0.98/0.002^2 = 4900$ puntos.
(b) $\sigma = 0.03$ y $n = 0.30 \times 0.70/0.03^2 = 233.3$, así que 234 puntos. (c) Para un error relativo
fijo, $n = (1-p)/(p\,\varepsilon^2)$ crece como $1/p$: el mineral accesorio exige 21 veces más
puntos, más de los que suele ofrecer una lámina.
\end{solution}

\begin{solution}{exo:b2:field-lab-methods-2:6}
(a) RGM-1: $72.89/73.45 = \qty{99.2}{\percent}$; BCR-2: $54.17/54.1 = \qty{100.1}{\percent}$.
(b) RGM-1: $0.19/\sqrt{20} = 0.042$, y $-0.56/0.042 = -13.2$; BCR-2: $0.13/\sqrt{20} = 0.029$, y
$0.07/0.029 = 2.4$. (c) RGM-1 tiene un sesgo: su media se aparta del valor recomendado 13 veces su propia
desviación típica, mucho más de lo que permite el azar (si el valor recomendado es correcto). BCR-2 está en
el límite de lo que da el azar. RGM-1 no es menos preciso: su desviación típica relativa,
$0.19/72.89 = \qty{0.26}{\percent}$, es próxima a la de BCR-2, \qty{0.24}{\percent}. El problema es
la veracidad, no la precisión.
\end{solution}

\begin{solution}{exo:b2:field-lab-methods-2:7}
(a) $M = 39.098 + 3 \times 24.305 + 26.982 + 3 \times 28.085 + 12 \times 15.999 + 2 \times 1.008 =
\qty{417.25}{g.mol^{-1}}$; los dos OH forman un \ce{H2O} ($18.015$), el \qty{4.32}{\percent}.
(b) \ce{K2O} $47.10/417.25 = \qty{11.3}{\percent}$; \ce{MgO} $3 \times 40.304/417.25 =
\qty{29.0}{\percent}$; \ce{Al2O3} $50.98/417.25 = \qty{12.2}{\percent}$; \ce{SiO2} $3 \times
60.083/417.25 = \qty{43.2}{\percent}$; total \qty{95.7}{\percent}. (c) La microsonda no
detecta el hidrógeno: a un análisis de mica le falta su agua, alrededor del \qty{4}{\percent}, y debería sumar cerca de
96.
\end{solution}

\begin{solution}{exo:b2:field-lab-methods-2:8}
(a) $\beta = \ln(0.1194/0.11905)/\ln(85.909261/87.905612) = 0.002936/(-0.02297) = -0.1278$.
(b) $0.71062 \times (86.908878/85.909261)^{-0.1278} = 0.70957$. (c) La 86/88 medida está por debajo de la
razón verdadera: se vio favorecido el isótopo pesado, y la 87/86 medida es demasiado alta.
\end{solution}

\begin{solution}{exo:b2:field-lab-methods-2:9}
(a) $0.1000 \times 15.271 \times 10^{-6}/83.9242 \times 1.3599 \times 10^{-3} = \qty{2.474e-11}{mol}$.
(b) Por \eqref{eq:b2:field-lab-methods-2:id}, $N_x = 2.474 \times 10^{-11} \times (732.8 -
0.874)/(0.874 - 0.0568) = \qty{2.216e-8}{mol}$ de $^{86}\mathrm{Sr}$; estroncio: $2.216 \times
10^{-8}/0.0986 \times 87.62 = \qty{1.969e-5}{g}$ en \qty{0.0500}{g}, \qty{394}{ppm} (393.9).
(c) Una recuperación del \qty{99.95}{\percent}: el método es exacto dentro del redondeo de la razón.
\end{solution}

\begin{solution}{exo:b2:field-lab-methods-2:10}
(a) $R_m = \sqrt{732.8 \times 0.05680} = 6.45$, donde $\mathcal M = 1.018$. (b) $\mathcal M = 2.31$
para 0.1, 1.13 para 0.5, 1.038 para 1.664 y 1.16 para 100. (c) Cerca de $R_x$ la mezcla apenas difiere de
la muestra, y cerca de $R_s$ del trazador, de modo que la diferencia que lleva la información se pierde en
el error de $R_m$: $\mathcal M$ diverge en ambos extremos. Entre ellos el mínimo es muy plano
($\mathcal M < 1.2$ desde alrededor de 0.4 hasta 100), así que el analista solo necesita el orden de magnitud correcto
del trazador, que proporciona una concentración aproximada, por ejemplo de XRF.
\end{solution}

\begin{solution}{exo:b2:field-lab-methods-2:11}
(a) $\sigma = \sqrt{0.62 \times 0.38/400} = 0.0243$: $62 \pm 2.4$. (b) $(65 - 62)/2.43 = 1.2$
desviaciones típicas: la roca bien puede ser una granodiorita. (c) $\sigma = 0.015$ requiere $0.62 \times
0.38/0.015^2 = 1047.1$, así que 1048 puntos de feldespato.
\end{solution}

\begin{solution}{exo:b2:field-lab-methods-2:12}
(a) Cada punto falla con probabilidad $1 - p$, independientemente: $(1-p)^n$. (b) $0.99^n \le 0.05$
da $n \ge \ln 0.05/\ln 0.99 = 298.1$: 299 puntos. (c) $n \ge \ln 0.05/\ln 0.9995 = 5990$: un recuento en malla
es la herramienta equivocada; un mineral accesorio escaso se busca recorriendo toda la lámina.
\end{solution}

\begin{solution}{pb:b2:field-lab-methods-2:1}
\noindent\textbf{Parte I.}
\textbf{1.} Cuarzo $372/1200 = \qty{31.0}{\percent}$, feldespato alcalino 29.0, plagioclasa 27.0, biotita
9.0, moscovita 3.0, óxidos opacos 1.0.
\textbf{2.} $\sqrt{0.31 \times 0.69/1200} = 0.0134$; moscovita: $\sqrt{0.97/(0.03 \times 1200)} =
\qty{16.4}{\percent}$.
\textbf{3.} $M = 9 + 3 + 1 = \qty{13.0}{\percent}$; $M' = 13 - 3 = \qty{10.0}{\percent}$ sin la
moscovita: leucocrática.
\textbf{4.} $Q + A + P = 1044$ puntos: Q $= 35.6$, A $= 33.3$, P $= 31.0$.
\textbf{5.} $324/672 = \qty{48.2}{\percent}$, con Q entre 20 y 60: un granito, y, entre 35 y
65, un monzogranito.
\textbf{6.} $\sqrt{0.482 \times 0.518/672} = 0.0193$: los límites de 35 y 65 quedan a 6.9 y 8.7
desviaciones típicas. El nombre es seguro.
\textbf{7.} $0.31 \times 0.69/0.01^2 = 2139$ puntos.

\noindent\textbf{Parte II.}
\textbf{8.} $0.0051/0.8000 = \qty{0.64}{\percent}$.
\textbf{9.} $0.7949/8.7949 = 0.0904$, y $72 \times 0.0904 = \qty{6.51}{\percent}$ de \ce{SiO2} en el
vidrio: la fusión diluye la muestra once veces.
\textbf{10.} $\bar I = 24.91$, $\bar c = 62.283$, $\sum(I - \bar I)^2 = 31.918$, $\sum(I - \bar I)(c -
\bar c) = 80.009$: $a = 2.5067$ por kcps, $b = 62.283 - 2.5067 \times 24.91 = -0.16$.
\textbf{11.} $2.5067 \times 28.63 - 0.16 = \qty{71.61}{\percent}$.
\textbf{12.} BCR-2 $-0.04$, AGV-2 $+0.05$, RGM-1 $-0.01$: todos menores que las desviaciones típicas
del método (de 0.1 a 0.2), así que la recta es adecuada en este intervalo.
\textbf{13.} $71.61 + 27.15 + 0.64 = \qty{99.40}{\percent}$, entre 97 y 101: no se repite.
\textbf{14.} $71.61 \times 100/(71.61 + 27.15) = \qty{72.51}{\percent}$.

\noindent\textbf{Parte III.}
\textbf{15.} $0.0986/0.8258 = 0.1194$: el valor de normalización es la razón natural.
\textbf{16.} $\beta = 0.007313/(-0.022972) = -0.3183$.
\textbf{17.} $0.71655 \times (86.908878/85.909261)^{-0.3183} = 0.71392$.
\textbf{18.} La 86/88 medida se aparta en $0.11853/0.1194 - 1 = \qty{-0.73}{\percent}$ del valor verdadero
en dos unidades de masa: alrededor del \qty{0.36}{\percent} por unidad de masa, a favor de los isótopos pesados.
\textbf{19.} La $^{86}\mathrm{Sr}/^{88}\mathrm{Sr}$ del trazador es 0.7661: la 86/88 de la mezcla ya no es
la natural, 0.1194, y fijarla en 0.1194 trataría el trazador como un sesgo instrumental.

\noindent\textbf{Parte IV.}
\textbf{20.} Suma de las razones a 88: $561.4 + 0.7661 + 0.1819 + 1 = 563.348$; $561.4/563.348 =
\qty{99.65}{\percent}$ y $0.7661/563.348 = \qty{0.13599}{\percent}$.
\textbf{21.} $0.2500 \times 15.271 \times 10^{-6}/83.9242 = \qty{4.549e-8}{mol}$ de estroncio, y
$\times 1.3599 \times 10^{-3} = \qty{6.186e-11}{mol}$ de $^{86}\mathrm{Sr}$.
\textbf{22.} $561.4/0.7661 = 732.8$ y $0.0056/0.0986 = 0.05680$.
\textbf{23.} $N_x = 6.186 \times 10^{-11} \times (732.8 - 1.6812)/(1.6812 - 0.0568) =
\qty{2.784e-8}{mol}$.
\textbf{24.} $\mathcal M = 1.6812 \times 732.74/(731.12 \times 1.6244) = 1.04$; la razón aporta
$1.04 \times \qty{0.01}{\percent}$, el trazador $0.031/15.271 = \qty{0.20}{\percent}$; en cuadratura,
\qty{0.20}{\percent}, dominado por la calibración del trazador.
\textbf{25.} Estroncio: $2.784 \times 10^{-8}/0.0986 = \qty{2.824e-7}{mol}$, es decir, $2.824 \times
10^{-7} \times 87.62 = \qty{2.474e-5}{g}$ en \qty{0.1000}{g}. El granito de Varden contiene
\textbf{\qty[text-series-to-math=true]{247}{ppm}} de estroncio, $\pm 0.5$.
\end{solution}
